\chapter{Dynamics of Rigid Body Systems}
\label{chap:3}
The purpose of a dynamics algorithm is to evaluate numerically an equation of motion; but, before we can evaluate it, we must first know what that equation is. This chapter introduces the equations of motion for a general rigid-body system, and presents a variety of methods for constructing them from simpler equations. These methods form the basis for various algorithms that appear in later chapters. However, the focus of this chapter is mainly on the mathematics of rigid-body dynamics, rather than on particular algorithms. A rigid-body system is a collection of rigid bodies that may be connected together by joints, and that may be acted upon by various forces. The effect of a joint is to impose a motion constraint on the two bodies it connects: relative motions are allowed in some directions but not in others. To be more precise, the effect of a joint is to introduce a constraint force into the system. This force has the special property that we do not know its value, but we do know what effect it has on the system—it takes whatever value it needs to take in order that the resulting motion complies with the motion constraint. A large part of this chapter is concerned with methods for describing motion constraints, and for applying them to rigid-body equations of motion. Sections 3.1 and 3.2 provide a high-level introduction to the subject: Section 3.1 presents the equations of motion in their various standard forms, and Section 3.2 explains how they are constructed from equations describing the subsystems, individual rigid bodies and motion constraints. The next three sections cover the topic of motion constraint in more detail: Section 3.3 shows how to express them as vector subspaces; Section 3.4 presents a standard classification of the kinds of constraint that can arise in a rigid-body system; and Section 3.5 presents a detailed model of joint constraints. (We define a joint to be any kinematic constraint on the relative motion of two bodies.) Finally, Sections 3.6 and 3.7 show how to obtain explicit equations of motion for individual

rigid bodies, and for collections of rigid bodies, that are connected together by joints. These last two sections illustrate the use of spatial vectors in formulating equations of motion. 3.1 Equations of Motion The equation of motion for a general rigid-body system can usually be written in the following canonical form: H(q) ¨q + C(q, ˙q) = τ . (3.1) In this equation, q, ˙q and ¨q are vectors of generalized position, velocity and acceleration variables, respectively, and τ is a vector of generalized forces. H is the generalized inertia matrix, and it is written H(q) to show that it depends on q. Likewise, C is the generalized bias force, and it is written C(q, ˙q) to show that it depends on both q and ˙q. Once these dependencies are understood, the shorter symbols H and C are used. The bias force is simply the value of τ that will produce zero acceleration. It accounts for the Coriolis and centrifugal forces, gravity, and any other forces acting on the system other than those in τ. Together, H and C are the coefficients of the equation of motion, and τ and ¨q are the variables. In addition to its role in the equation of motion, H has another important property: the kinetic energy of the system is T = 1 2 ˙qTH ˙q . (3.2) Every quantity in Eq. 3.1 is either a coordinate vector or a matrix. It is useful to consider, at least briefly, what kinds of object these quantities represent. To this end, we let the symbols Mn and Fn denote the spaces of n-dimensional motion and force vectors, respectively, where n is the degree of freedom of the system described by Eq. 3.1. We can now say that ˙q and ¨q represent elements of Mn, whereas τ and C represent elements of Fn, and H is a mapping from Mn to Fn. However, q does not represent a vector, but instead represents a point in the system’s configuration space, which is an n-dimensional manifold. By definition, a set of generalized force variables cannot be chosen at random, but must be chosen so that τ T ˙q is the power delivered by τ to the system. This implies that a system of generalized coordinates on Mn and Fn is actually a system of dual coordinates. This, in turn, implies that generalized motion and force vectors share many of the algebraic properties of spatial vectors, as described in Chapter 2. The one big exception is that a cross-product operator is not defined on generalized motions and forces. Strictly speaking, H and C do not depend only on q and ˙q, but also on the rigid-body system itself. Thus, it would be more accurate to write H(model, q) and C(model, q, ˙q), where the symbol model refers to a collection of data that describes a particular rigid-body system in terms of its component parts: the

number of bodies and joints, the manner in which they are connected together, the values of the bodies’ inertia parameters, and so on. We call this kind of description a system model, and it is the subject of Chapter 4. For now, it is enough to know that such a model exists, and that quantities like H and C depend on it. Our eventual aim is to be able to calculate the numeric values of one or more of the quantities appearing in Eq. 3.1. In particular, we have an interest in the following two calculations: forward dynamics, which is the calculation of ¨q given τ, and inverse dynamics, which is the calculation of τ given ¨q. We can express these calculations succinctly as follows: ¨q = FD(model, q, ˙q, τ) (3.3) and τ = ID(model, q, ˙q, ¨q) , (3.4) where FD and ID are functions denoting the forward and inverse dynamics calculations, respectively. On comparing these equations with Eq. 3.1, it is clear that FD and ID must evaluate to H−1(τ −C) and H¨q + C, respectively. However, the algorithms that implement them need not necessarily work by calculating C or H or H−1. Thus, the useful property of Eqs. 3.3 and 3.4 is that they show clearly the inputs and outputs of each calculation, but do not imply any particular method of calculation. Equation 3.1 assumes that the position variables are the integrals of the velocity variables. This is not always the case. For example, if the system contains nonholonomic constraints (see §3.4) then it will have more position variables than velocity variables. Alternatively, certain kinds of motion are better described by a redundant set of position variables, which again leads to there being more position variables than velocity variables. Even if the number of position and velocity variables is the same, it may occasionally be useful for the latter to be something other than the derivatives of the former. Whatever the reason, we can adapt Eq. 3.1 simply by replacing ˙q with an alternative vector of velocity variables, α, to get H(q) ˙α + C(q, α) = τ . (3.5) A corresponding modification is then made to Eqs. 3.3 and 3.4. For simulation purposes, we still need to know the value of ˙q, so that it can be integrated to obtain q. Equation 3.5 must therefore be supplemented with an equation to calculate ˙q from q and α. This equation can be written ˙q = Q(q) α , (3.6) where Q is a matrix that depends on q and the system model. Alternatively, it can be written ˙q = qdfn(model, q, α) , (3.7)

where qdfn denotes the calculation of ˙q from q and α. Equation 3.6 reveals the mathematical structure of the relationship between q, ˙q and α by showing that ˙q depends linearly on α; while Eq. 3.7 shows the inputs and outputs of the calculation without implying any particular calculation method. Still on the subject of simulation, most simulators expect to be given a set of first-order differential equations. We can accommodate this by defining a state variable, x = [qT αT]T, and expressing the forward dynamics in state-space form: ˙x =  ˙q ˙α  =  Q α H−1(τ −C)  . (3.8) If FDx denotes the state-space forward dynamics calculation function, then ˙x = FDx(model, x, τ) =  qdfn(model, q, α) FD(model, q, α, τ)  . (3.9) In subsequent sections and chapters, we will work mostly with systems that can be described by Eq. 3.1, rather than the more general systems described by Eq. 3.5. This is because most of the results we shall obtain for the simpler kind of system can be adapted in a straightforward manner to the more general kind. Most dynamics algorithms are indifferent to the distinction between Eqs. 3.1 and 3.5, except at the level of programming details. Finally, there is one last possibility to mention. It is sometimes impractical or undesirable to identify and work with an explicit set of independent velocity variables. In this case, one still uses Eq. 3.1 or 3.5 as the equation of motion; but now ˙q and ¨q (or α and ˙α) are subject to a set of motion constraints, and so the equation of motion must be supplemented with a second equation to describe these constraints. This topic is covered in the next section. 3.2 Constructing Equations of Motion The equations of motion for a rigid-body system are obtained as the end result of a sequence of mathematical operations. The starting point is a collection of equations of motion for individual rigid bodies and/or subsystems; and the two main operations that we perform on them are: • collecting equations together to form the equation of a bigger subsystem, and • applying additional motion constraints. By performing these operations in a particular order, we define a particular procedure for obtaining the desired equation of motion, hence also a particular algorithm (or class of algorithms) for a computer to calculate the numeric value of that equation. Some of the most common procedures are as follows. 1. Collect all of the bodies (i.e., collect their equations into a single equation), and then apply all of the constraints.

2. Obtain the equation of motion for the spanning tree, and then apply all remaining constraints. 3. Starting with a single body or subsystem do: add one more body or subsystem; apply constraints pertaining to that body or subsystem; repeat. 4. Combine subsystems in pairs; apply constraints separately within each subsystem; repeat. Method 1 is typical of many simple dynamics algorithms. It is easy to understand, but it produces large, sparse matrices. Algorithms that take this approach must use sparse matrix techniques, or they will be too slow compared with other algorithms. An example of this method appears in Section 3.7. Method 2 is standard for closed-loop systems, and is covered in detail in Chapter 8. The spanning tree of a closed-loop system is defined in Chapter 4 to be a kinematic tree that accounts for all of the bodies in the original system but only a subset of the motion constraints. The point of method 2 is that there are very efficient algorithms available for calculating the dynamics of the spanning tree. Method 3 is approximately how the articulated-body algorithm works, except that it uses articulated-body equations instead of complete equations of motion. This algorithm is described in Chapter 7. Finally, method 4 is used by divide-and-conquer algorithms (Featherstone, 1999a,b). Algorithms that work like this are suitable for implementation on parallel computers. Collecting Equations of Motion Suppose we have a set of N rigid-body subsystems, such that system i has the equation of motion Hi¨qi +Ci = τi. If we wish to treat them as a single system, then the equation of motion for that system is

H1 0 · · · 0 0 H2 · · · 0 ... ... ... ... 0 0 · · · HN

¨q1 ¨q2 ... ¨qN

+

C1 C2 ... CN

=

τ1 τ2 ... τN

. (3.10) Note that a single rigid body is also a rigid-body subsystem. Thus, we could have started with a set of N rigid bodies, or even an arbitrary mixture of individual bodies and other systems. If subsystem i happens to be a single rigid body, then we can use its spatial equation of motion, Iiai + pi = fi, in place of the generalized-coordinates equation of motion (i.e., we use Ii in place of Hi, etc.). The reason why this works is because spatial motion and force vectors are special cases of generalized motion and force vectors; that is, M6 and F6 are special cases of Mn and Fn, and Pl¨ucker coordinates on M6 and F6 are a special case of generalized coordinates on Mn and Fn. Thus, we can mix spatial velocities freely with generalized velocities, and similarly for accelerations and forces.

Motion Constraints Motion constraints are algebraic constraints on the motion variables of a system. They can be expressed in two different ways, as follows: position velocity acceleration implicit: φ(q) = 0 K ˙q = 0 K¨q = k explicit: q = γ(y) ˙q = G ˙y ¨q = G¨y + g (3.11) where K = ∂φ ∂q , k = −˙K ˙q , G = ∂γ ∂y and g = ˙G ˙y . (These are not the most general equations. The general case is covered in §3.4.) If φ and γ both describe the same constraint, then φ ◦γ = 0 , KG = 0 and Kg = k . (3.12) Motion constraints are caused by constraint forces. If an unconstrained (or less constrained) rigid-body system has the equation of motion H¨q + C = τ , (3.13) and this system is to be subjected to a motion constraint, then the equation of motion for the constrained system is H¨q + C = τ + τc , (3.14) where τc is the constraint force. This force is unknown, but it has one important property that allows us either to calculate its value, or to eliminate it from the equation, as desired: The constraint force delivers zero power along every direction of velocity freedom that is compatible with the motion constraints. (Jourdain’s principle of virtual power.) In effect, this statement says that τc must satisfy τc · ˙q = 0, not just for one value of ˙q, but for every value that is allowed by the constraint. If the system is subject to an implicit motion constraint, then it can be shown that τc takes the form τc = KTλ , (3.15) where λ is a vector of unknown force variables; and if the system is subject to an explicit motion constraint, then τc has the property GTτc = 0 . (3.16) These equations are easily verified. If τc = KTλ then τc · ˙q = λTK ˙q, and this expression will be zero for every ˙q satisfying K ˙q = 0. Likewise, if ˙q = G ˙y then ˙q · τc = ˙yTGTτc, and this expression will be zero for all ˙y if GTτc = 0. The elements of λ can be regarded as a set of Lagrange multipliers.

Applying Implicit Motion Constraints To apply an implicit motion constraint to a rigid-body system, we simply substitute Eq. 3.15 into Eq. 3.14, and combine the result with the implicit acceleration constraint equation in Eq. 3.11. The result is  H KT K 0   ¨q −λ  = τ −C k  . (3.17) This is a set of n + nc equations in n + nc unknowns, where n = dim(¨q) is the degree of freedom of the unconstrained system and nc = dim(λ) = dim(φ) is the number of constraints imposed upon it. It can be shown that the rank of the coefficient matrix is n + nic, where nic = rank(K) is the number of independent constraints imposed on the system. Thus, if nic = nc then the coefficient matrix is nonsingular, and Eq. 3.17 can be solved for both ¨q and λ. However, if nic < nc then it is still possible to solve Eq. 3.17 for ¨q, but there will be nc −nic degrees of indeterminacy in λ. A system in which nic < nc is said to be overconstrained. Ideally, we would always have nic = nc, but it is not always possible to guarantee this in practice. Applying Explicit Motion Constraints To apply an explicit motion constraint to a rigid-body system, we combine Eq. 3.14 with Eq. 3.16 and the explicit acceleration constraint equation in Eq. 3.11. The resulting equation is

H −1 0 −1 0 G 0 GT 0

¨q τc ¨y

=

τ −C −g 0

. (3.18) Applying one step of Gaussian elimination to this equation produces

H −1 0 0 H−1 G 0 GT 0

¨q τc ¨y

=

τ −C H−1(τ −C) −g 0

, from which we can extract H−1 G GT 0  τc ¨y  =  H−1(τ −C) −g 0  ; (3.19) and applying Gaussian elimination to this equation produces H−1 G 0 −GTHG   τc ¨y  =  H−1(τ −C) −g −GT(τ −C −Hg)  , from which we get GTHG ¨y = GT(τ −C −Hg) . (3.20)

Finally, if we make the substitution u = GTτ, where u is a new set of generalized forces, then we get HG ¨y + CG = u , (3.21) where HG = GTHG and CG = GT(C + Hg) . Thus, we have obtained a new equation of motion for the constrained system that has the same algebraic form as the equation of motion for the unconstrained system in Eq. 3.13. A few remarks are in order at this point. First, Eq. 3.19 can be regarded as the dual of Eq. 3.17 in the sense that the motion and force variables have changed places. However, 3.17 is more useful than 3.19 in practice. Second, Eq. 3.20 can be obtained directly from Eq. 3.14 by first premultiplying it by GT, which eliminates τc, and then substituting for ¨q using the explicit acceleration constraint in Eq. 3.11. This method is sometimes called the projection method on the grounds that the first step amounts to a projection of the equation of motion onto the range space of GT. Third, HG inherits from H the properties of symmetry and positive definiteness. It also has the property that the kinetic energy of the constrained system is T = 1 2 ˙yTHG ˙y . (3.22) Finally, we can show that u· ˙y is the power delivered to the constrained system as follows: u · ˙y = (GTτ)T ˙y = τ TG ˙y = τ · ˙q . (3.23) Thus, u does indeed satisfy the definition of a vector of generalized forces. 3.3 Vector Subspaces Vector subspaces are the natural mathematical tool for describing and analysing kinematic constraints because they capture the essential aspects of a constraint. For example, the system in Eq. 3.13 is free to move in Mn because ˙q is an element of Mn and it is unconstrained. However, the system in Eq. 3.14 is constrained to move only in the subspace S ⊂Mn. If the constraint has been described implicitly, then S = null(K); and if it has been described explicitly, then S = range(G). If K1 and K2 both describe the same constraint, then the one thing they must have in common is that null(K1) = null(K2) = S. Likewise, if G1 and G2 describe the same constraint, then the one thing they must have in common is that range(G1) = range(G2) = S. Thus, it is the subspace S, rather than any one particular matrix, that captures the essence of the constraint. This section therefore reviews some basic material on vector subspaces and how they can be used.

Basics A subspace of a vector space is a subset that is also a vector space. Thus, a subspace is any subset that is closed with respect to the operations of addition and scalar multiplication. Let V be a vector space, and let Y = \{y1, y2, . . .\} be any subset of V . The quantity span(Y ) is the set of all linear combinations of the elements of Y , and is therefore a subspace. It follows that any subset that satisfies Y = span(Y ) is a subspace. If S is a subspace of V , and Y = \{y1, y2, . . .\} is a basis on S, then S = span(Y ) and dim(S) = |Y | (i.e., the dimension of S equals the number of elements in Y ). A zero-dimensional subspace is a set containing only the zero vector. Given that a subspace is itself a vector space, it follows that a subspace may have subspaces, and so on. There is no special symbol for ‘is a subspace of’. Instead, we use the symbol for a subset. Thus, we write Y ⊆V to indicate that Y is a subset of V , and S ⊆V to indicate that S is a subspace. One must read the surrounding text to discover which meaning is intended. The expression S ⊂V indicates that S is a proper subspace of V , which implies that dim(S) < dim(V ). Suppose that S ⊆V , and that dim(S) = m and dim(V ) = n. It follows that m must lie in the range 0 ≤m ≤n. If m = 0 then S = \{0\}; and if m = n then S = V ; but if m takes any other value, then there are infinitely many subspaces having that dimension. The number of parameters required to identify uniquely an m-dimensional subspace of an n-dimensional vector space is m(n −m). Matrix Representation The simplest way to represent a subspace is as the range or null space of a matrix. We will mostly use the former. Let S be an m-dimensional subspace of the n-dimensional vector space V , and let S = \{s1, . . . , sm\} be a basis on S. Any vector v ∈S can then be expressed in the form v = Pm i=1 si αi, where αi are the coordinates of v in S. If si are themselves coordinate vectors, expressed in a basis on V , then they can be assembled into an n × m matrix, S = [s1 · · · sm], which has the property that range(S) = S. v can then be expressed in the form v = Sα, where α = [α1 · · · αm]T. The matrix S defines both a subspace and a basis. In fact, of the mn numbers contained in S, we can say that m(n −m) define the subspace and m2 define the basis. If S′ = \{s′ 1, . . . , s′ m\} is another basis on S, then the matrix S′ = [s′ 1 · · · s′ m] defines the same subspace as S, but a different basis. In this case, S′ can be expressed in the form S′ = SA, where A is the coordinate transform from S′ to S coordinates. (The elements of A satisfy s′ j = Pm i=1 si Aij.) Thus, any matrix SA, where A is invertible, describes the same subspace as S. Decomposition of Vectors One of the most basic operations that can be performed on a vector is to decompose it into its components in two subspaces. Given a vector v ∈V and

subspaces S1, S2 ⊆V , the idea is to express v in the form v = v1 + v2, where v1 ∈S1 and v2 ∈S2. This decomposition is possible if v ∈span(S1 ∪S2), and it is unique if S1 ∩S2 = \{0\} (i.e., S1 and S2 have no nonzero element in common). If we want the decomposition to be both possible and unique for any v, then S1 and S2 must satisfy S1 ∩S2 = \{0\} and dim(S1) + dim(S2) = dim(V ) . (3.24) (These conditions imply span(S1 ∪S2) = V .) If S1 and S2 satisfy Eq. 3.24 then we can say that V is the direct sum of S1 and S2, which is written V = S1 ⊕S2 . (3.25) This equation is simply a short-hand way of saying that any element of V can be decomposed uniquely into the sum of an element of S1 and an element of S2. Essentially, Eqs. 3.24 and 3.25 say the same thing. Let us now perform the decomposition. We are given a vector v ∈V and two matrices S1 and S2 representing subspaces that satisfy Eq. 3.25; and the task is to find α1 and α2 that satisfy v = S1 α1 + S2 α2 = [S1 S2] α1 α2  . (3.26) This problem can be solved using the following result: if S1 ⊕S2 = V then [S1 S2] defines a basis on V , and is therefore a nonsingular matrix. Given that [S1 S2]−1 exists, the solution to Eq. 3.26 is simply  α1 α2  = [S1 S2]−1v . (3.27) This formula extends in the obvious way to the task of decomposing a vector into components in more than two subspaces. Orthogonal Complements If two Euclidean vectors, u and v, satisfy the equation u · v = 0, then they are said to be orthogonal. If Y1 and Y2 are subsets of En, and every element of Y1 is orthogonal to every element of Y2, then the two subsets are said to be orthogonal, which is written Y1 ⊥Y2. The set of all vectors that are orthogonal to a given subset, Y ⊆En, is called its orthogonal complement, and is written Y ⊥. An orthogonal complement is always a subspace. Some basic properties of orthogonal subsets are listed in Table 3.1. This form of orthogonality is based on the Euclidean inner product, and is therefore applicable only to Euclidean vectors. However, it is possible to define other forms of orthogonality, which are based on other scalar products. In particular, if u ∈U and v ∈V , where U = V ∗, then a scalar product is defined between u and v, and it is possible to say that if u · v = 0 then they

Subsets: Subspaces: Y1 ⊥Y2 ⇒Y2 ⊥Y1 dim(S) + dim(S⊥) = dim(U) Y1 ⊥Y2 ⇒span(Y1) ⊥Y2 S1 ⊥S2 ⇔ST 1 S2 = 0 Y1 ⊥Y2 ∧Z ⊆Y1 ⇒Z ⊥Y2 S1 ⊕S⊥ 2 = U ⇔(ST 1 S2)−1 exists Y1 ⊥Y2 ⇒Y1 ⊆Y ⊥ 2 S1 ⊕S⊥ 2 = U ⇔S⊥ 1 ⊕S2 = V Y ⊥= (span(Y ))⊥ Euclidean Subspaces: (Y ⊥)⊥= span(Y ) S ⊕S⊥= En Table 3.1: Properties of orthogonal subsets. In the general case, S, S1, Y, Y1 ⊆U and S2, Y2 ⊆V where U = V ∗. In the Euclidean case, S, . . . , Y2 ⊆En are orthogonal. By extension, if Y1 and Y2 are subsets of U and V , respectively, and every element of Y1 is orthogonal to every element of Y2, then Y1 ⊥Y2. Likewise, the set of all elements of V that are orthogonal to Y ⊆U is the orthogonal complement of Y . Observe that the Euclidean and dual forms of orthogonality have different and incompatible mathematical properties. In particular, the relationship Y1 ⊥Y2 can only be true in the Euclidean sense if Y1 and Y2 are subsets of the same vector space, and it can only be true in the dual sense if they are subsets of different vector spaces.1 Orthogonal complements have gained a degree of notoriety in the robotics literature, owing to the mistaken practice of treating 6D vectors as if they were Euclidean, and using the Euclidean form of orthogonal complement instead of the dual form. The matter is discussed in Duffy (1990). Some authors have sought to distance themselves from this mistake by using alternative names like ‘natural orthogonal complement’ or ‘reciprocal complement’ for the dual form of the orthogonal complement. Orthogonal complements play the following role in rigid-body dynamics: if a system is initially free to move in Mn, and a motion constraint restricts the motion to a subspace S ⊆Mn, then the constraint force that imposes this constraint on the system is an element of S⊥⊆Fn. Example 3.1 Suppose we have been given a motion subspace, S ⊆M6, and a spatial inertia, I : M6 7→F6, which is positive definite. With these quantities, we can define two subspaces, Ta = IS and Tc = S⊥, which have the property Ta ⊕Tc = F6. (The expression IS is the image of S under the mapping I, i.e., the set \{Is | s ∈S\}.) We can prove this property as follows. According to Eq. 3.24, we must show that dim(Ta) + dim(Tc) = 6, and that Ta ∩Tc does not contain a nonzero element. The first part is easy: dim(Ta) = dim(S) and dim(Tc) = 6 −dim(S). To prove the second part, we suppose that t̸ = 0 ∈ (Ta ∩Tc) and prove a contradiction: if t ∈Ta then there exists a nonzero s ∈S such that t = Is; but if t ∈Tc then sTt = 0, implying sTIs = 0; but this is 1One way around this is to declare that a Euclidean vector space is self-dual.

impossible because I is positive definite. Having established that Ta⊕Tc = F6, we can decompose any f ∈F6 uniquely into components fa ∈Ta and fc ∈Tc. Given any matrix S spanning S, we can obtain fa and fc as follows: STf = STfa + STfc = STfa , but fa = ISα for some α, so STf = STISα , which implies α = (STIS)−1STf , hence fa = IS(STIS)−1STf . (3.28) fc then follows simply from fc = f −fa. The expression in Eq. 3.28 clearly must be independent of the particular choice of S. We can prove this by replacing each instance of S in this equation with SA, where A is any invertible matrix of the correct size, and showing that the ‘A’s cancel out. The physical interpretation of this decomposition is as follows. If I is the inertia of a rigid body that is constrained to move in the subspace S, and f is a force applied to that body, then fa is the component of f that causes acceleration, and fc is the component that is opposed by the constraint force (i.e., the constraint force will be −fc). If the body is at rest, and is not subject to any forces other than f and the motion constraint force, then its acceleration will be a = I−1fa = S(STIS)−1STf . The expression S(STIS)−1ST is the apparent inverse inertia of the constrained rigid body (cf. Eq. 3.54). 3.4 Classification of Constraints Kinematic constraints can be classified according to the nature of the constraint they impose on a rigid-body system. In every case, the constraint is described by an algebraic equation in the motion variables of the system, but the form of the equation varies. A standard classification hierarchy is shown in Figure 3.1. The first distinction is made between equality and inequality constraints. The former arise from permanent physical contact between pairs of rigid bodies; and the latter arise when bodies are free both to make contact and to separate. Inequality constraints can be used to describe phenomena such as collision, bouncing and loss of contact. They are the subject of Chapter 11, and will not be considered further here. The next distinction is between holonomic and nonholonomic constraints. The former are constraints on the position variables, and arise typically from

kinematic constraints inequality equality nonholonomic scleronomic holonomic rheonomic φ(...)=0 φ(q1,...,qn,t)=0 φ(...) 0 φ(q1,...,qn)=0 φ(q1,...,qn,q1,...,qn,t)=0 . . Figure 3.1: Classification of kinematic constraints sliding contact. The latter are constraints on the velocity variables, and arise typically from rolling contact. A nonholonomic constraint function, φnh, is linear in the velocity variables, and can therefore be expressed in the form φnh(q, ˙q, t) = φ0 nh + n X i=1 φi nh ˙qi , (3.29) where the coefficients φi nh are functions of the position variables and time only. For comparison, the time derivative of a holonomic constraint function, φh, is d dtφh(q, t) = ∂φh ∂t + n X i=1 ∂φh ∂qi ˙qi , (3.30) which has the same algebraic form. Thus, at the velocity and acceleration levels, there is no difference between holonomic and nonholonomic constraints. The big difference between these two equations is that Eq. 3.29 is nonintegrable; that is, φnh is not the derivative of any function, or else it would simply be a holonomic constraint expressed via its derivative.2 The consequence of this property is that a system containing nonholonomic constraints has more positional degrees of freedom than velocity degrees of freedom, and therefore needs more position variables than velocity variables. If a rigid-body system contains at least one nonholonomic constraint, then it is called a nonholonomic system; otherwise, it is called a holonomic system. Thus, Eq. 3.1 can describe only a holonomic system, but Eq. 3.5 can describe both holonomic and nonholonomic systems. The final distinction is between a scleronomic constraint, which is a function of the position variables only, and a rheonomic constraint, which is a function 2Given that a constraint equation holds at all times, rather than only at a single instant, the two equations ˙φ = 0 and φ = 0 describe the same constraint.

B1 B2 B1 B2 l1 l2 l3 (a) (b) (c) (d) Figure 3.2: Examples of kinematic constraints also of time. Scleronomic constraints arise from sliding contact between a surface fixed in one body and a surface fixed in another. They are both the simplest and the most common kind of constraint in a typical mechanical system. Rheonomic constraints can be regarded as scleronomic constraints in which one or more sliding freedoms have been forced to vary as a prescribed function of time. The role of rheonomic constraints is to introduce kinematic excitations into a system; that is, prescribed motions. Some examples of kinematic constraints are shown in Figure 3.2. Diagrams (a) and (b) show a cylindrical joint and a prismatic joint, respectively. Both are scleronomic constraints. The cylindrical joint allows body B2 to slide and rotate relative to B1, whereas the prismatic joint allows only a relative sliding motion. Thus, the cylindrical joint allows two degrees of motion freedom, while the prismatic joint allows only one. If we were to make one of the cylindrical joint’s freedoms a prescribed function of time, then it would become a one-degree-offreedom rheonomic constraint. If we were to do the same to the prismatic joint, then it would become a zero-degree-of-freedom rheonomic constraint (which is actually quite useful). A kinematic constraint is nearly always a relationship between two bodies. However, it is possible to devise constraints that involve more than two bodies, and that cannot be reduced to a set of two-body constraints. An example is shown in Figure 3.2(c). This diagram shows a 2-D rigid-body system in which three identical circular bodies are constrained to lie within the perimeter of a massless, zero-diameter, inelastic string. Assuming the string stays taut, the three bodies are free to move in any way that satisfies the equation l1 + l2 + l3 + 2πr = p, where p and r are the perimeter of the string and the radius of a body, respectively. Figure 3.2(d) shows a sphere resting on a planar surface. If this sphere is constrained to roll without slipping, then it has three degrees of instantaneous motion freedom: it can roll in two directions, and it can spin about the contact point. However, by a suitable sequence of manoeuvres, it is possible for the sphere to arrive at any point on the plane, with any orientation. It therefore

has five degrees of position freedom. This is the characteristic feature of a nonholonomic constraint. If the sphere were free to slide, then it would be a scleronomic constraint. It is clear from the shapes of the bodies in Figure 3.2(a) and (b) that they cannot be pulled apart; but there is nothing that physically prevents the string in 3.2(c) from going slack, or the sphere in 3.2(d) from losing contact with the plane. If we know enough about the forces that will be acting on these systems to be sure that such an event will not occur, then it is possible to model these constraints as equality constraints; otherwise, they must be modelled as inequality constraints. 3.5 Joint Constraints We define a joint to be a kinematic constraint between two bodies. Thus, a joint can impose anywhere from zero to six constraints on the relative velocity of a pair of rigid bodies. We call the two bodies the predecessor and the successor, and we say that the joint connects from the predecessor to the successor. The joint velocity, vJ, is then defined to be the velocity of the successor relative to the predecessor: vJ = vs −vp , (3.31) where vp and vs are the velocities of the predecessor and successor bodies, respectively. The general form of the joint constraint is vJ = S(q, t) ˙q + σ(q, t) , (3.32) where S is the joint’s motion subspace matrix, σ is the bias velocity, q and ˙q are the joint’s position and velocity variables, and t is time. If the position variables are not the integrals of the velocity variables, then one must use a different symbol either for ˙q or for q. The bias velocity is the value that vJ takes when ˙q = 0, and it has a nonzero value only if the joint depends explicitly on time (i.e., if it is a rheonomic or time-dependent nonholonomic constraint). If the joint does not depend on time, then Eq. 3.32 simplifies to vJ = S(q) ˙q . (3.33) The effect of Eq. 3.32 is to confine the velocity vJ −σ to the subspace S ⊆M6. If the joint allows nf degrees of relative motion freedom, then dim(S) = nf and S is a 6 × nf matrix.3 The number of constraints imposed by the joint is nc = 6 −nf, and the constraint force must lie in the nc-dimensional subspace S⊥⊆F6. 3If nf = 0 then S is a 6×0 matrix. Such quantities follow the basic rules of matrix algebra, plus one extra rule: the product of an m × 0 matrix with a 0 × n matrix is an m × n matrix of zeros. Thus, the product of the 6 × 0 matrix S with the 0 × 1 vector ˙q is a 6 × 1 vector of zeros.

Let fJ be the force transmitted across the joint. We define this to be the force transmitted from the predecessor to the successor; so the force acting on the successor is fJ and the force acting on the predecessor is −fJ. This force can be regarded as the sum of an active force and a constraint force. The former comes from actuators, springs, dampers, or any other force element acting at the joint (or between the two bodies), and the latter is the force that imposes the motion constraint. At this point, we have a choice. The constraint force must be an element of the subspace S⊥, but there are two ways to handle the active force: it can be treated either as a known spatial force, or as a generalized force mapped onto a subspace of F6. We shall choose the latter here. Let T and Ta be the constraint and active force subspaces, respectively. By definition, T = S⊥, but Ta can be any subspace that satisfies T ⊕Ta = F6. We can now express fJ in the form fJ = Taτ + T λ , (3.34) where Taτ is the active force, T λ is the constraint force, Ta and T satisfy T T a S = 1 (3.35) and T TS = 0 , (3.36) τ is the generalized force vector corresponding to ˙q, and λ is a vector of constraint force variables. Given the properties of Ta and T , we can immediately deduce that STfJ = τ (3.37) and T TvJ = T Tσ . (3.38) Equation 3.37 is the formula for obtaining the generalized force at a joint from the spatial force transmitted across it; and Eq. 3.38 is the implicit form of the velocity constraint equation. The justification for Eq. 3.35 is as follows. The total power delivered by the joint is fJ · vJ, but this comes from two independent power sources: σ and τ. The power attributable to σ is fJ · σ, so the power attributable to τ must be fJ · S ˙q. Now, this second power must also be equal to τ · ˙q because they are the generalized force and velocity vectors at the joint, so we have the following power balance equation: τ · ˙q = fJ · S ˙q = τ TT T a S ˙q ; (3.39) but this equation must be true for all τ and ˙q, so we must have T T a S = 1.

The joint’s acceleration constraint equations are obtained directly from Eqs. 3.32 and 3.38 by differentiation. The derivative of Eq. 3.32 is aJ = S¨q + ˙S ˙q + ˙σ = S¨q + (˚ S + vs × S) ˙q + (˚σ + vs × σ) = S¨q + ˚ S ˙q + ˚σ + vs × (S ˙q + σ) = S¨q + cJ + vs × vJ , (3.40) where ˚ S and ˚σ are the apparent derivatives of S and σ in a coordinate system that is moving with the successor body, and cJ = ˚ S ˙q + ˚σ . (3.41) (Actually, cJ = ˚vJ.) If S and σ are functions of q and t, as shown in Eq. 3.32, then the general formulae for ˚ S and ˚σ are ˚ S = ∂S ∂t + np X i=1 ∂S ∂qi ˙qi (3.42) and ˚σ = ∂σ ∂t + np X i=1 ∂σ ∂qi ˙qi , (3.43) where np is the number of joint position variables. Before moving on, readers should bear in mind that the relatively complicated equations above apply to a general case that is almost never needed in practice. This is partly because most common joint types have the property that ˚ S = 0 and σ = 0 (implying cJ = 0), and partly because on the rare occasions when one encounters an explicit time dependency in a joint (i.e., a joint for which σ̸ = 0), it is almost always possible to implement the time-dependent term via the hybrid dynamics algorithms of Chapter 9 rather than using σ. Finally, the implicit acceleration constraint equation can be obtained either by differentiating Eq. 3.38 or by multiplying Eq. 3.40 by T T. Both give the same result, but the latter is easier: T TaJ = T T(S¨q + cJ + vs × vJ) = T T(cJ + vs × vJ) . (3.44) Example 3.2 Let us examine how to define an active force subspace. Figure 3.3 shows a system comprising a moving body, B, that is connected to a fixed base via a revolute joint, and a motor, M, that drives this joint via a pinion and gear wheel. The x and y axes of a coordinate frame are also shown, the z axis

y x r f M B Figure 3.3: Defining an active force subspace being aligned with the joint’s rotation axis. Expressed in these coordinates, the joint’s motion and constraint-force subspace matrices are S =

0 0 1 0 0 0

and T =

1 0 0 0 0 0 1 0 0 0 0 0 0 0 0 0 0 1 0 0 0 0 0 1 0 0 0 0 0 1

. Strictly speaking, these are not the only possible values for S and T , since any two matrices that span the correct subspaces would do. The active force is transmitted to B at the point where the two gears touch, and it is a linear force of magnitude f in the y direction. The coordinates of the contact point are (−r, 0, 0), where r is the radius of the gear wheel, so the Pl¨ucker coordinates of the active spatial force are f = [0 0 −rf 0 f 0]T. Now, τ is just a scalar in this example, so if f = Taτ then Ta is just a scalar multiple of f. From Eq. 3.35 we have that Ta must satisfy T T a S = 1, so the solution is simply Ta = f/(f TS) =

0 0 1 0 −1/r 0

. Observe that both Ta and range(Ta) depend on r. This dependency has no effect on the equation of motion, but it does affect the value of the constraint force. Physically, the constraint force comes from the joint’s bearings. Thus, if the task is to simulate the motion of the system, then r is irrelevant; but if the task is to select appropriate bearings for the joint, then the value of r must be taken into account.

f fc I v a S Figure 3.4: A constrained rigid body 3.6 Dynamics of a Constrained Rigid Body Consider the rigid body shown in Figure 3.4. This body is connected to a fixed base by a joint having a motion subspace of S ⊆M6. The joint imposes nc constraints on the motion of the body, leaving it with nf degrees of motion freedom, where nf = dim(S) and nc = 6 −nf = dim(S⊥). Three forces act on this body: an applied force, f, a constraint force, fc, and a gravity force, fg, which is not shown. The equation of motion for this body is therefore f + fc + fg = Ia + v ×∗Iv , where I, a and v are the body’s inertia, acceleration and velocity, respectively. We are not particularly interested in the individual terms fg and v ×∗Iv, so we collect them into a single bias-force term, p = v ×∗Iv −fg. The equation of motion for the constrained rigid body is then f + fc = Ia + p , (3.45) subject to the constraints v ∈S and fc ∈S⊥. (3.46) We seek a solution to this system of equations in which the acceleration is given as a function of the applied force. Note that fc is an unknown quantity, and must be evaluated and/or eliminated as part of the solution process. Method 1 Let us introduce a 6 × nf matrix, S, that spans the subspace S. With this matrix, the constraint equations in Eq. 3.46 can be written in the form v = S ˙q (3.47) and STfc = 0 ; (3.48)

and the acceleration constraint is a = S ¨q + ˙S ˙q , (3.49) where ˙q and ¨q are vectors of joint velocity and acceleration variables, respectively. Equation 3.48 implies that we can eliminate fc from Eq. 3.45 by multiplying both sides by ST, giving STf = ST(I a + p) . (3.50) Substituting Eq. 3.49 into this equation produces STf = ST(I (S ¨q + ˙S ˙q) + p) , which can be solved for ¨q as follows: ¨q = (STIS)−1ST(f −I ˙S ˙q −p) . (3.51) The expression STIS is an nf × nf symmetric, positive-definite matrix, and is therefore guaranteed to be invertible. Having found an expression for ¨q, the final step is to substitute Eq. 3.51 back into Eq. 3.49, giving a = S (STIS)−1ST(f −I ˙S ˙q −p) + ˙S ˙q . (3.52) This equation can be written in the form a = Φ f + b , (3.53) where Φ = S (STIS)−1ST (3.54) and b = ˙S ˙q −Φ(I ˙S ˙q + p) . (3.55) Φ is the constrained body’s apparent inverse inertia, and b is its bias acceleration, which is the acceleration it would have if f = 0. The apparent inverse inertia describes how the acceleration would change in response to a change in the applied force. It is a symmetric, positive-semidefinite matrix with the following properties: range(Φ) = S, null(Φ) = S⊥and rank(Φ) = nf. Method 2 Another way to solve the system of equations in Eqs. 3.45 and 3.46 is to introduce a 6 × nc matrix, T , that spans S⊥. With this matrix, the constraint equations can be written in the form T Tv = 0 (3.56) and fc = T λ , (3.57)

where λ is a vector of constraint force variables, and the acceleration constraint is T Ta + ˙T Tv = 0 . (3.58) Substituting Eq. 3.57 into Eq. 3.45 gives f + T λ = I a + p . (3.59) Equations 3.58 and 3.59 can then be combined into a single equation,  I T T T 0   a −λ  = f −p −˙T Tv  , (3.60) which can be solved for a and λ. The coefficient matrix in this equation is symmetric and nonsingular, but not positive definite. Method 3 A third possibility is to derive an equation of motion expressed in generalized coordinates. The vectors ˙q and ¨q can be used as the generalized velocity and acceleration variables of the constrained rigid body, so Eq. 3.51 already gives the generalized acceleration as a function of the applied spatial force, f. Let τ denote the generalized force acting on the body. For τ to be equivalent to f, both forces must deliver the same power for any velocity in S. The power delivered by τ is τ T ˙q, and that delivered by f is f Tv, so we have τ T ˙q = f Tv = f TS ˙q ; but this equation must hold for all ˙q, so τ = STf . (3.61) Substituting this equation into Eq. 3.51 gives ¨q = (STIS)−1(τ −ST(I ˙S ˙q + p)) , which is the equation of motion for the constrained rigid body, expressed in generalized coordinates. It can be written in the standard form H ¨q + C = τ , (3.62) where H = STIS (3.63) is the generalized inertia matrix and C = ST(I ˙S ˙q + p) (3.64) is the generalized bias force.

Equation 3.61 implies a small loss of information in going from f to τ. Infinitely many different values of f map to the same value of τ, and therefore produce the same acceleration, but each produces a different value of fc. If only τ is known, then it is possible to work out the value of the sum f + fc, but not the values of the individual vectors. If the objective is motion simulation, then the lost information is irrelevant; but if the objective is to design a mechanical system, then the lost information is relevant to tasks like calculating the dynamic load forces on the joint bearings. 3.7 Dynamics of a Multibody System We conclude this chapter by developing an equation of motion for a system containing multiple bodies and joints. This analysis shows the connection between the joint-related material in Sections 3.5 and 3.6 and the generalized constraints in Section 3.2, by showing how the latter can be constructed from the former. Suppose we have a rigid-body system containing the following elements: a fixed base, which we regard as body number zero; NB mobile rigid bodies, which are numbered 1 to NB; and NJ joints, which are numbered 1 to NJ. For each joint j, we define p(j) and s(j) to be the body numbers of the predecessor and successor bodies of that joint. This set of 2NJ numbers defines the connectivity of the system. For the analysis to follow, we make no restrictions on the connectivity. In formulating the equation of motion, we shall follow Method 2 in Section 3.6, but with extra steps where necessary because of the greater complexity of the current system. Step 1: Body Equations of Motion We assume that the only forces acting on body i are the force of gravity, fgi, and forces from the joints that are connected to it. The equation of motion for body i can then be written fi = Ii ai + pi , where fi is the sum of the joint forces acting on body i, and pi = vi×∗Iivi−fgi is the bias force. The equations for all NB bodies can be assembled into a single equation,

f1 f2 ... fNB

=

I1 0 · · · 0 0 I2 · · · 0 ... ... ... ... 0 0 · · · INB

a1 a2 ... aNB

+

p1 p2 ... pNB

. (3.65) We can express this equation more compactly as f = I a + p , (3.66)

where f, a and p are 6NB-dimensional vectors, and I is a 6NB × 6NB blockdiagonal matrix. For completeness, we also define v = [vT 1 vT 2 · · · vT NB]T, where vi is the velocity of body i, as this vector will be needed in the joint constraint equation. Step 2: Joint Motion from Body Motion The vector a in Eq. 3.66 is a vector of body accelerations; but the constraints imposed by the joints are expressed in terms of joint velocities and accelerations. Therefore, we need a formula to express joint motion in terms of body motion. Let vJj and aJj denote the velocity and acceleration, respectively, across joint j; and let vJ and aJ be the 6NJ-dimensional vectors formed by concatenating vJ1 · · · vJNJ and aJ1 · · · aJNJ, respectively. By definition, vJj and aJj are the velocity and acceleration of the joint’s successor body relative to its predecessor; so vJj = vs(j) −vp(j) and aJj = as(j) −ap(j) . These equations can be collected together and expressed as follows: vJ =

vJ1 ... vJNJ

=

vs(1) −vp(1) ... vs(NJ) −vp(NJ)

= P Tv (3.67) and aJ =

aJ1 ... aJNJ

=

as(1) −ap(1) ... as(NJ) −ap(NJ)

= P Ta , (3.68) where Pij =

16×6 if i = s(j) −16×6 if i = p(j) 06×6 otherwise. (3.69) P is therefore a 6NB × 6NJ matrix, which is organized as an NB × NJ block matrix of 6 × 6 blocks, and it maps the vectors of collected body velocities and accelerations (v and a) to the vectors of collected joint velocities and accelerations (vJ and aJ). Each row4 in P corresponds to a body, and each column to a joint. (Observe that we are using i as a body-number index and j as a joint-number index.) Each column contains at most two nonzero blocks: a 1 in the row corresponding to the joint’s successor body and a −1 in the row corresponding to its predecessor. If a joint connects between two moving bodies, then it will have exactly two nonzero blocks in its column, as just described. 4Strictly speaking, we mean block row, and similarly for columns.

However, if it connects to or from the fixed base (body 0), then its column will contain only one nonzero block: a 1 or −1, as appropriate, in the row for the moving body it connects to/from. Each row i will contain a 1 for each joint that has body i as its successor, and a −1 for each joint that has body i as its predecessor. Step 3: Body Forces from Joint Forces Let fJj be the force transmitted across joint j, and let fJ be the 6NJ-dimensional vector formed by concatenating the vectors fJ1 · · · fJNJ. By definition, fJj is the force transmitted from body p(j) to body s(j) through joint j. If we were to define the sets s−1(i) and p−1(i) to be the sets of all joints having body i as their successor or predecessor, respectively, then the vector fi (which was defined earlier to be the sum of all joint forces acting on body i) could be expressed in the following manner: fi = X j∈s−1(i) fJj − X j∈p−1(i) fJj . However, on comparing this expression with the description of row i of P , it is clear that we could also express fi as follows: fi = NJ X j=1 PijfJj . Therefore f = P fJ . (3.70) Step 4: Motion Constraints Let Tj be the 6 × ncj matrix that spans the subspace S⊥ j , where Sj is the motion subspace for joint j and ncj is the number of constraints it imposes. The velocity constraint equation for joint j can then be written T T j vJj = 0 , and the acceleration constraint is T T j aJj + ˙T T j vJj = 0 . Collecting together the acceleration constraints gives T TaJ + ˙T TvJ = 0 , (3.71) where T is the 6NJ × nc block-diagonal matrix having Tj in its jth diagonal block. nc is the total number of constraints imposed by the joints, and is given by nc = PNJ j=1 ncj. Combining Eq. 3.71 with Eqs. 3.67 and 3.68 yields the following constraint on the body accelerations: T TP Ta + ˙T TP Tv = 0 . (3.72)

Step 5: Constraint Forces fJj is the sum of an active force and a constraint force, and it can be expressed in the form fJj = Tajτj + Tjλj (see Eq. 3.34). In this equation, Taj is the 6×nfj matrix that spans the activeforce subspace for joint j, τj is the nfj-dimensional vector of generalized forces at joint j, λj is the ncj-dimensional vector of constraint-force variables, and nfj = 6 −ncj is the number of motion freedoms allowed by joint j. Collecting the equations for every joint, we have fJ = Taτ + T λ , (3.73) where τ and λ are the vectors formed by concatenating τ1 · · · τNJ and λ1 · · · λNJ, respectively, and Ta is the block-diagonal matrix having Taj on its jth diagonal block. Step 6: Final Equation Substituting Eqs. 3.70 and 3.73 into 3.66 produces the following equation of motion: P (Taτ + T λ) = I a + p ; and combining this with Eq. 3.72 produces  I P T T TP T 0   a −λ  = P Taτ −p −˙T TP Tv  . (3.74) This is the equation of motion for the original rigid-body system. The two unknowns are a and λ, and the equation can be solved uniquely for a. If the coefficient matrix is nonsingular, then λ is also unique. On comparing the equations here with those in Section 3.2, it can be seen that T TP T and −˙T TP Tv here are equivalent to K and k in Eq. 3.11, that I, a and p here are equivalent to H, ¨q and C in Eq. 3.14, that P Taτ and P T λ are equivalent to τ and τc in Eq. 3.14, and that Eq. 3.74 as a whole is equivalent to Eq. 3.17. Discussion Equation 3.74 is a viable basis for a dynamics algorithm, but there are several details that would need to be resolved in the process of designing the algorithm. For example, the velocity and acceleration variables are clearly the Pl¨ucker coordinates of the velocity and acceleration vectors of the individual bodies, but in what coordinate system? No mention has been made of the position variables, or of how to calculate I, T , Taτ or ˙T Tv. Another issue is how to solve Eq. 3.74. The coefficient matrix is large, but it is also sparse, and using a sparse factorization procedure would greatly improve the efficiency. However, if the constraints in Eq. 3.72 are not all linearly independent, then the coefficient

matrix will be singular, in which case it will be necessary to use a solution procedure that works with singular matrices. Finally, there is the problem of constraint stabilization: as Eq. 3.72 specifies the motion constraints at the acceleration level, it is possible for errors to accumulate in the position and velocity constraints, due to numerical rounding and truncation errors during simulation. This problem is discussed in Section 8.3.
